\documentclass[12pt]{article}
\usepackage[T1]{fontenc}
\usepackage[utf8]{inputenc}
\usepackage{lmodern}
\usepackage{textcomp}
\usepackage{enumitem}
\usepackage{geometry}
\usepackage{graphicx}
\usepackage{amsmath, amssymb}
\geometry{margin=1in}
\begin{document}
\begin{center}
Physics - Graduate Level Assessment 1 \
Total Marks: 40 \
Answer all questions.
\end{center}

\section*{Multiple Choice (10 marks)}
\begin{enumerate}[label=\arabic*.]
\item 10 calories of heat are supplied to 1 g of water. How much does the mass of the water increase?
\begin{enumerate}[label=(\alph*)]
    \item 5.5 \ensuremath{\times} 10^-13 g
    \item 2.5 \ensuremath{\times} 10-13 g
    \item 2.0 \ensuremath{\times} 10^-12 g
    \item 4.7 \ensuremath{\times} 10-13 g
    \item 8.6 \ensuremath{\times} 10^-13 g
\end{enumerate}
\item Compute the mean molecular speed v in the light gas hydrogen (H$_{2}$) in m/s
\begin{enumerate}[label=(\alph*)]
    \item 1750.0
    \item 1500.0
    \item 2200.0
    \item 2100.0
    \item 1400.0
\end{enumerate}
\item Consider two mercury barometers, one having a cross-sectional area of 1 cm 2 and the other 2 cm 2. The height of mercury in the narrower tube is
\begin{enumerate}[label=(\alph*)]
    \item two-thirds
    \item twice
    \item the same
    \item half
    \item four times higher
\end{enumerate}
\item Which one of these constellations is not located along the Milky Way in the sky?
\begin{enumerate}[label=(\alph*)]
    \item Virgo
    \item Perseus
    \item Ursa Major
    \item Sagittarius
    \item Aquila
\end{enumerate}
\item An object, starting from rest, is given an acceleration of 16 feet per $second^2$ for 3 seconds. What is its speed at the end of 3 seconds?
\begin{enumerate}[label=(\alph*)]
    \item 48 ft per sec
    \item 64 ft per sec
    \item 30 ft per sec
    \item 72 ft per sec
    \item 16 ft per sec
\end{enumerate}
\end{enumerate}

\section*{True / False (10 marks)}
\textit{Answer True or False}
\begin{enumerate}[label=\arabic*.]
\setcounter{enumi}{5}
\item True or False: The angular magnification of the refracting telescope can be calculated as 20, based on the separation of the lenses.
\item True or False: The rest mass of a particle with total energy 5.0 GeV and momentum 4.9 GeV/c is approximately 1.0 GeV/$c^2$.
\item True or False: The efficiency of a device that converts 100 J of input energy into 40 J of useful work is 25\%.
\item True or False: Doubling the thickness of a neutral density filter will always halve its transmittance regardless of initial transmission values.
\item True or False: If the volume of an object doubles while its mass remains constant, its density will halve.
\end{enumerate}

\section*{Long Form Response (20 marks)}
\textit{Answer each question with a clear and complete explanation.}
\begin{enumerate}[label=\arabic*.]
\setcounter{enumi}{10}
\item Discuss the shape of the illuminated area on a screen positioned 1 cm and 2 m away from a tiny plane mirror measuring 2 mm \ensuremath{\times} 2 mm, explaining the underlying physical principles.
\item Considering a converging lens with a 5.0 cm focal length used as a simple magnifier, determine the object distance required to produce a virtual image 25 cm from the eye. Additionally, calculate the magnification and discuss the implications of these values in the context of optical systems.
\end{enumerate}

\vfill
\noindent\textit{End of Paper}
\end{document}